\documentclass[10pt,a4paper]{amsart}
\usepackage{amsmath,amssymb,amsfonts,amsthm}
\usepackage{enumitem}
\usepackage[left=1.5cm,right=1.5cm]{geometry}

\newcommand{\R}{\mathbb{R}}
\renewcommand{\d}{\mathrm{d}}
\newcommand{\Lap}{\mathcal{L}}

\theoremstyle{definition}
\newtheorem{problema}{Problema}

\title{Introducao as Equacoes Diferenciais -- Shortlist para a Prova 2}

\begin{document}
\maketitle

\noindent{\small Problemas selecionados a partir de \texttt{Prova2-Sugestoes.tex}: P1, P5, P7, P8, P9, P10, P12, P13, P15, P17, P20, P22.}

\vspace{4mm}

\begin{problema}[P1]
\textbf{[Facil; Lista 4, equacoes lineares de segunda ordem.]}
Resolva
\[
    y''-2y'+y=e^t(1+t),
    \qquad y(0)=0,\qquad y'(0)=1.
\]
\end{problema}

\begin{problema}[P5]
\textbf{[Facil; Lista 4, exponencial de matriz.]}
Calcule $e^{tA}$ e resolva $x'=Ax$ para
\[
    A=\begin{pmatrix}
        2&1&0\\
        0&2&0\\
        0&0&-1
    \end{pmatrix}.
\]
\end{problema}

\begin{problema}[P7]
\textbf{[Facil; Lista 4, transformada inversa.]}
Encontre a transformada inversa de Laplace.
\begin{enumerate}[label=(\alph*)]
    \item $F(s)=\dfrac{2s+3}{s^2+4s+13}$.
    \item $F(s)=\dfrac{e^{-3s}}{s-2}$.
    \item $F(s)=\dfrac{1+e^{-2s}}{s^2+6}$.
\end{enumerate}
\end{problema}

\begin{problema}[P8]
\textbf{[Facil; Lista 4, equacao integral vira PVI.]}
Seja $f:[0,\infty)\to\R$ continua e suponha que
\[
    f(x)=1+\int_0^x (x-t)\bigl(f(t)-2e^t\bigr)\,\d t
    \qquad(x\ge 0).
\]
Mostre que $f$ satisfaz um problema de valor inicial de segunda ordem e
resolva-o explicitamente.
\end{problema}

\begin{problema}[P9]
\textbf{[Medio-dificil; Lista 4, ideia central: comparar derivadas/series.]}
Sejam $A,B$ matrizes reais ou complexas. Prove que
\[
    e^{t(A+B)}=e^{tA}e^{tB}
\]
para todo $t\in\R$ se e somente se
\[
    AB=BA.
\]
\end{problema}

\begin{problema}[P10]
\textbf{[Medio-dificil; Lista 4, ideia central: somar iteracoes integrais.]}
Seja $f_0:[0,1]\to\R$ uma funcao continua. Para $n\ge 1$, defina
\[
    f_n(x)=\int_0^x f_{n-1}(t)\,\d t.
\]
Prove que a serie
\[
    \sum_{n=1}^{\infty} f_n(x)
\]
converge para todo $x\in[0,1]$, e encontre uma formula explicita para sua
soma.
\end{problema}

\begin{problema}[P12]
\textbf{[Medio-dificil; Lista 4, sistemas matriciais.]}
Considere matrizes $A,B\in\R^{2\times2}$ e a equacao diferencial matricial
\[
    X'(t)=AX(t)+X(t)B,\qquad X(t)\in\R^{2\times2}.
\]
\begin{enumerate}[label=(\alph*)]
    \item Prove que $X(t)=e^{At}X(0)e^{Bt}$ e solucao.
    \item Se
    \[
        A=\begin{pmatrix}0&-\omega\\ \omega&0\end{pmatrix},
        \qquad
        B=\begin{pmatrix}\alpha&0\\0&\beta\end{pmatrix},
    \]
    escreva explicitamente a solucao $X(t)$ para
    $X(0)=\begin{pmatrix}x_{11}&x_{12}\\x_{21}&x_{22}\end{pmatrix}$.
    \item Prove que, caso $B=-A$, entao, para todo $k\in\mathbb{N}$, o traco
    $\operatorname{Tr}(X(t)^k)$ e constante ao longo do tempo.
\end{enumerate}
\end{problema}

\begin{problema}[P13]
\textbf{[Dificil; Lista 4, Jacobi--Liouville.]}
Seja $A:[0,T]\to M_n(\R)$ continua e seja $X:[0,T]\to M_n(\R)$ a solucao de
\[
    X'(t)=A(t)X(t),
    \qquad X(0)=I.
\]
\begin{enumerate}[label=(\alph*)]
    \item Prove a identidade de Jacobi
    \[
        \frac{\d}{\d t}\det X=\det X\cdot\operatorname{tr}(X^{-1}X')
    \]
    enquanto $X$ for invertivel.
    \item Prove a formula de Liouville
    \[
        \det X(t)=\exp\left(\int_0^t \operatorname{tr}(A(s))\,\d s\right),
    \]
    e conclua, em particular, que $X(t)$ e invertivel para todo $t$.
\end{enumerate}
\end{problema}

\begin{problema}[P15]
\textbf{[Dificil; Lista 4, estimativa assintotica.]}
Seja $f(x)$ uma funcao real tal que
\[
    \lim_{x\to+\infty}\frac{f(x)}{e^x}=1
\]
e
\[
    |f''(x)|<c|f'(x)|
\]
para todo $x$ suficientemente grande. Prove que
\[
    \lim_{x\to+\infty}\frac{f'(x)}{e^x}=1.
\]
\end{problema}

\begin{problema}[P17]
\textbf{[Dificil; banco/Lista 4, series e funcoes geradoras.]}
Seja $B_0=1$ e
\[
    B_{n+1}=\sum_{k=0}^n \binom{n}{k}B_k.
\]
Defina
\[
    B(t)=\sum_{n\ge 0}B_n\frac{t^n}{n!}.
\]
Deduza uma equacao diferencial para $B(t)$, resolva-a e, expandindo a
resposta em serie de potencias, deduza a formula de Dobinski
\[
    B_n=e^{-1}\sum_{m=0}^{\infty}\frac{m^n}{m!}.
\]
\end{problema}

\begin{problema}[P20]
\textbf{[Ultra-dificil / candidato a impossivel; Lista 4, pendulo com
comprimento variavel.]}
Sejam $l_0,c,\alpha,g$ constantes positivas, e seja $x(t)$ a solucao da
equacao diferencial
\[
    \bigl((l_0+ct^\alpha)^2x'\bigr)'
    +g[l_0+ct^\alpha]\sin x=0,
    \qquad t\ge 0,\qquad -\frac{\pi}{2}<x<\frac{\pi}{2},
\]
satisfazendo as condicoes iniciais
\[
    x(t_0)=x_0,\qquad x'(t_0)=0.
\]
Esta e a equacao do pendulo matematico cujo comprimento varia de acordo com
a lei $l=l_0+ct^\alpha$. Prove que $x(t)$ esta definida no intervalo
$[t_0,\infty)$; alem disso, se $\alpha>2$, entao, para todo $x_0\neq 0$,
existe $t_0$ tal que
\[
    \liminf_{t\to\infty}|x(t)|>0.
\]
\end{problema}

\begin{problema}[P22]
\textbf{[Ultra-dificil; Lista 4, equacao integral com retardo.]}
Seja $x:[0,\infty)\to\R$ uma funcao diferenciavel que satisfaz
\[
    x'(t)
    =
    -2x(t)\sin^2 t
    +
    (2-|\cos t|+\cos t)
    \int_{t-1}^{t}x(s)\sin^2 s\,\d s
\]
em $[1,\infty)$. Prove que $x$ e limitada em $[0,\infty)$ e que
\[
    \lim_{t\to\infty}x(t)=0.
\]
A conclusao continua verdadeira para funcoes que satisfazem
\[
    x'(t)
    =
    -2x(t)
    +
    (2-|\cos t|+\cos t)
    \int_{t-1}^{t}x(s)\,\d s\,?
\]
\end{problema}

\end{document}
